\subsection{Equivalence Relations}
\label{subsec:equivalence-relations}

\begin{tcolorbox}[title={Definition --- Equivalence Relation}]
\begin{definition}[Equivalence Relation]
\label{def:equiv-relation}
A binary relation $R$ on $A$ is an \textbf{equivalence relation} if it is
reflexive, symmetric, and transitive.
We write $a \sim b$ when $R$ is an equivalence relation and the particular
relation is understood from context.
\end{definition}
\end{tcolorbox}

\begin{remark*}[Standard quantified statement]
\[
\begin{aligned}
\operatorname{EquivalenceRelation}(R,A)
\Longleftrightarrow\;&
R\subseteq A\times A \\
&{}\land
[(\forall a \in A)\, a \mathrel{R} a] \\
&{}\land[(\forall a,b \in A)\, a \mathrel{R} b \Rightarrow b \mathrel{R} a] \\
&{}\land[(\forall a,b,c \in A)\,
a \mathrel{R} b \land b \mathrel{R} c \Rightarrow a \mathrel{R} c].
\end{aligned}
\]
\end{remark*}

\begin{remark*}[Definition predicate reading]
$R$ is an equivalence relation on $A$ exactly when it is reflexive, symmetric,
and transitive.
\end{remark*}

\begin{remark*}[Negated quantified statement]
\[
\neg\operatorname{EquivalenceRelation}(R,A)
\Longleftrightarrow
\neg\operatorname{ReflexiveRelation}(R,A)
\lor
\neg\operatorname{SymmetricRelation}(R,A)
\lor
\neg\operatorname{TransitiveRelation}(R,A).
\]
\end{remark*}

\begin{remark*}[Interpretation]
An equivalence relation on $A$ is precisely the data needed to partition $A$
into disjoint, exhaustive equivalence classes.  Quotient groups, quotient
rings, quotient vector spaces, and quotient topologies are all built by
choosing equivalence relations compatible with the relevant operations.
\end{remark*}

\begin{remark*}[Dependencies]
\begin{itemize}
  \item Binary relation.
  \item Reflexive relation.
  \item Symmetric relation.
  \item Transitive relation.
\end{itemize}
\end{remark*}
